\documentclass[a4paper,12pt,twocolumn]{article}

\usepackage[landscape,top=1cm,bottom=1.4cm,left=1cm,right=1cm,columnsep=1cm]{geometry}
\usepackage[fleqn]{amsmath}
\usepackage{amssymb}
\usepackage{fontspec}
\usepackage{enumitem}
\usepackage{xeCJK}
\usepackage[hidelinks]{hyperref}

\setlength{\mathindent}{1.5em}

% Set fonts
\setmainfont{Times New Roman}
\setsansfont{Arial}
\setmonofont{JetBrains Mono}

% Chinese font support
\setCJKmainfont{Iansui}[
  BoldFont=Iansui,
  AutoFakeBold=3.17
]

\pagestyle{plain}
\setlength{\columnseprule}{0.2pt}

\newcommand{\examdate}{2026/03/18}
\newlength{\problemnumberwidth}
\settowidth{\problemnumberwidth}{\textbf{62.}\ }
\newcommand{\problemtitle}[2]{%
  \hypertarget{#1}{}%
  \par\vspace{0.4em}%
  \noindent{\large\bfseries #2}\par
  \vspace{0.3em}%
  \hrule
  \vspace{0.9em}%
}
\newcommand{\worksheettoc}{%
  \noindent{\Large\bfseries Contents}\par
  \vspace{0.3em}%
  \hrule
  \vspace{0.9em}%
  \begin{tabular}{@{}p{0.26\columnwidth}p{0.66\columnwidth}@{}}
    \hyperlink{sec:10-1}{\textbf{10-1}} & \hyperlink{sec:10-1}{16, 19, 22, 57} \\
    \hyperlink{sec:10-2}{\textbf{10-2}} & \hyperlink{sec:10-2}{42} \\
    \hyperlink{sec:10-3}{\textbf{10-3}} & \hyperlink{sec:10-3}{57, 61, 62} \\
    \hyperlink{sec:10-4}{\textbf{10-4}} & \hyperlink{sec:10-4}{44} \\
  \end{tabular}\par
}
\newcommand{\problem}[2]{%
  \par\noindent
  \hangindent=\problemnumberwidth
  \hangafter=1
  \makebox[\problemnumberwidth][l]{\textbf{#1.}}#2\par
  \vspace{1.4em}%
}
\newlist{subparts}{enumerate}{1}
\setlist[subparts]{%
  label=(\alph*),
  labelindent=\problemnumberwidth,
  leftmargin=*,
  labelsep=0.5em,
  align=left,
  itemsep=0.2em,
  topsep=0.3em,
  parsep=0pt,
  partopsep=0pt
}

\begin{document}

\twocolumn[
{\centering
\Large\bfseries Calculus Problems\par
\vspace{0.5em}
\noindent\hfill\normalsize\normalfont \examdate\par
\vspace{0.8em}
}
]

\worksheettoc
\vfill\null
\newpage

\problemtitle{sec:10-1}{10-1}

\problem{16}{%
\begin{subparts}
\item Eliminate the parameter to find a Cartesian equation of the curve.
\item Sketch the curve and indicate with an arrow the direction in which the curve is traced as the parameter increases.
\end{subparts}
\[
x=\csc t,\ y=\cot t,\ 0<t<\pi
\]}

\problem{19}{%
\begin{subparts}
\item Eliminate the parameter to find a Cartesian equation of the curve.
\item Sketch the curve and indicate with an arrow the direction in which the curve is traced as the parameter increases.
\end{subparts}
\[
x=\ln t,\ y=\sqrt{t},\ t\geq 1
\]}

\problem{22}{%
\begin{subparts}
\item Eliminate the parameter to find a Cartesian equation of the curve.
\item Sketch the curve and indicate with an arrow the direction in which the curve is traced as the parameter increases.
\end{subparts}
\[
x=\sinh t,\ y=\cosh t
\]}

\problem{57}{Find the point at which the parametric curve intersects itself and the corresponding values of \(t\).
\begin{subparts}
\item \(x=1-t^2,\ y=t-t^3\)
\item \(x=2t-t^3,\ y=t-t^2\)
\end{subparts}}

\problemtitle{sec:10-2}{10-2}

\problem{42}{Let \(\mathcal{R}\) be the region enclosed by the loop of the curve in Example 1.
\begin{subparts}
\item Find the area of \(\mathcal{R}\).
\item If \(\mathcal{R}\) is rotated about the \(x\)-axis, find the volume of the resulting solid.
\item Find the centroid of \(\mathcal{R}\).
\end{subparts}
\medskip
\noindent\textit{Hint:} Use the parametric curve from Example 1:
\[
x=t^2,\ y=t^3-3t
\]
You can first find the intersection point(s) of the loop and the corresponding parameter values (for example, \(t=\pm\sqrt{3}\) gives the point \((3,0)\)), then use parametric integrals for area, volume about the \(x\)-axis, and centroid.}

\problemtitle{sec:10-3}{10-3}

\problem{57}{Show that the polar equation
\[
r=a\sin\theta+b\cos\theta,\qquad ab\ne 0
\]
represents a circle. Find its center and radius.}

\problem{61}{Graph the polar curve. Choose a parameter interval that produces the entire curve.
\[
r=e^{\sin\theta}-2\cos(4\theta)
\]
\noindent(butterfly curve)}

\problem{62}{Graph the polar curve. Choose a parameter interval that produces the entire curve.
\[
r=\lvert\tan\theta\rvert^{\lvert\cot\theta\rvert}
\]
\noindent(valentine curve)}

\problemtitle{sec:10-4}{10-4}

\problem{44}{Find the area of the shaded region bounded by
\[
r^2=\sqrt{3}\sin 2\theta
\]
and
\[
r=\sqrt{2}\cos\theta
\]}

\end{document}
